\documentclass[border=12pt]{standalone}
\usepackage{ctex}
\usepackage{amsmath,amssymb}
\usepackage{xcolor}
\setlength{\parindent}{0pt}
\setlength{\parskip}{6pt}

\definecolor{titlegreen}{RGB}{14,90,69}
\definecolor{linegreen}{RGB}{42,151,121}
\definecolor{bggreen}{RGB}{245,252,250}

\newcommand{\stage}[2]{%
  \vspace{2mm}
  \noindent\fcolorbox{linegreen}{bggreen}{%
    \begin{minipage}{0.95\linewidth}
      \textbf{#1}\\[2pt]
      #2
    \end{minipage}%
  }%
}

\begin{document}
\begin{minipage}{22cm}
{\LARGE\bfseries\color{titlegreen} LLM 场景：一次 REINFORCE 训练迭代}

\textbf{场景：}输入问题“什么是机器学习？”，模型逐 token 生成回答。

\stage{阶段 A：采样回答}{
假设模型生成：\texttt{机器学习是一种让计算机从数据中学习的技术。}

状态/动作示例：
\[
\begin{aligned}
&s_0=\text{“什么是机器学习？”}\\
&a_0=\text{“机器”},\quad s_1=\text{“什么是机器学习？机器”}\\
&a_1=\text{“学习”},\quad s_2=\text{“什么是机器学习？机器学习”}
\end{aligned}
\]
}

\stage{阶段 B：获取回报}{
把完整回答送入奖励模型，得到序列级奖励 $R=0.8$。

在“仅终局给奖励”且 $\gamma=1$ 的简化设定下：
\[
G_0=G_1=\cdots=0.8
\]
}

\stage{阶段 C：计算梯度并更新}{
\[
\nabla_\theta\log\pi_\theta(\text{“机器”}\mid s_0)\cdot G_0,
\qquad
\nabla_\theta\log\pi_\theta(\text{“学习”}\mid s_1)\cdot G_1
\]
把所有时间步梯度求和后更新参数。
}

\stage{阶段 D：效果}{
正回报会增强高质量回答 token 的概率；负回报会抑制低质量回答 token 的概率。
}
\end{minipage}
\end{document}
